\documentclass[11pt]{article}

% =====================================================
% Fonts and math
% =====================================================
\usepackage{mathpazo}   % Palatino text + math
\usepackage{amsmath, amssymb}

% =====================================================
% Page layout
% =====================================================
\usepackage{fullpage}
\usepackage{parskip}

% =====================================================
% Lists
% =====================================================
\usepackage{enumitem}

\begin{document}

\begin{center}
    {\Large \textbf{Lecture 7 In-Classs Worksheet: Should a Firm Raise Its Price?}}
\end{center}

\vspace{1em}

\section*{Setup}

A firm wants to \textbf{maximize profit}. It currently charges a single price $P$ for its good.

Assume:
\begin{itemize}
    \item The firm charges \textbf{one price} (no price discrimination).
    \item \textbf{Marginal cost is positive:} $MC > 0$.
    \item The firm considers raising its price \textbf{a little} (a small increase in $P$).
\end{itemize}

Recall:
\[
    \pi = TR - TC
\]
\[
    TR = P \cdot Q(P)
\]

\vspace{1em}

\section*{Part A — Warm-up: What changes when price rises?}

When the firm raises price slightly:

\begin{enumerate}[label=\arabic*.]
    \item Quantity demanded $Q$ \textbf{(rises / falls)}.
    \item Total revenue $TR$ may \textbf{(rise / fall)} depending on demand elasticity.
    \item Total cost $TC$ changes because output changes and $MC > 0$.
\end{enumerate}

\textbf{Key idea:} If the firm sells fewer units, it saves production costs on those units.

\vspace{1em}

\section*{Part B — Elastic Demand}

Suppose demand at the current price is \textbf{elastic}.

\subsection*{Question 1: Revenue}

If demand is elastic, a small increase in price causes quantity to fall by a \textbf{larger percentage} than the price rises.

\begin{itemize}
    \item Does $TR = P \cdot Q$ increase or decrease?

\end{itemize}

Explain briefly:
\vspace{2em}

\subsection*{Question 2: Cost}

When price rises, quantity falls, so the firm produces fewer units.

\begin{itemize}
    \item Does total cost $TC$ increase or decrease?
\end{itemize}

Why does the assumption $MC > 0$ matter here?

\vspace{3em}

\subsection*{Question 3: Profit}

Profit is given by:
\[
    \pi = TR - TC
\]

Given your answers above, when demand is elastic and the firm raises price a little, what happens to profit $\pi$?
Explain briefly:
\vspace{3em}

\newpage

\section*{Part C — Inelastic Demand}

Suppose demand at the current price is \textbf{inelastic}.

\subsection*{Question 4: Revenue}

If demand is inelastic, a small increase in price causes quantity to fall by a \textbf{smaller percentage} than the price rises.

\begin{itemize}
    \item Does $TR = P \cdot Q$ increase or decrease?

\end{itemize}

Explain briefly:
\vspace{2em}

\subsection*{Question 5: Cost}

When price rises, quantity falls, so the firm produces fewer units.

\begin{itemize}
    \item Does total cost $TC$ increase or decrease?
          \[
              \text{$TC$ \hspace{1cm} (increases \ / \ decreases)}
          \]
\end{itemize}

\vspace{1em}

\subsection*{Question 6: Profit}

Profit is given by:
\[
    \pi = TR - TC
\]

Given your answers above, when demand is inelastic and the firm raises price a little:

\[
    \pi \hspace{1cm} (\text{increases} \ / \ \text{decreases} \ / \ \text{ambiguous})
\]

Explain briefly:
\vspace{3em}

\section*{Part D — Synthesis}

\subsection*{Fill in the rule}

\begin{itemize}
    \item If demand is \textbf{elastic}, raising price a little tends to \underline{\hspace{4cm}} profit.
    \item If demand is \textbf{inelastic}, raising price a little tends to \underline{\hspace{4cm}} profit.
\end{itemize}

\subsection*{One-sentence intuition}

Complete the sentence:

\begin{quote}
    Raising price is more attractive when demand is \underline{\hspace{3cm}} because revenue moves \underline{\hspace{3cm}}, and costs move \underline{\hspace{3cm}} when output falls (since $MC > 0$).
\end{quote}

\vspace{1em}


\end{document}